% ============================================================
% 1. Introduction
% ============================================================
\section{Introduction}
\label{sec:introduction}

\subsection{Problem Statement}
\label{sec:problem_statement}

Entity resolution (ER)—the task of identifying records that refer to the same real-world entity across one or more data sources—is a foundational challenge in data integration.
The problem arises from pervasive data quality issues: attribute ambiguity (e.g., `J. Smith` vs. `John Smith`), data entry errors, missing values, temporal drift in attributes, and schema heterogeneity.
Traditional deterministic approaches (exact matching, rule-based systems) are labour-intensive, brittle, and fail to generalise across diverse corruption patterns.

ER encompasses a wide spectrum of problem structures and inference methods: clustering-based, probabilistic models, supervised learning, etc. An annex summarising these methodologies is provided in \Cref{annex:er_methodologies}.

The choice of method is dictated not by preference but by the \textit{constraints} of the specific ER problem at hand.
This report addresses ER under the following \textbf{explicit constraints}:

\begin{enumerate}
	\item \textbf{Global entity resolution}: The goal is to resolve entities across two datasets (Source A and Source B) into a unified global view.
	\item \textbf{Unsupervised learning}: No ground-truth labels are available for training; parameter estimation must be data-driven.
	\item \textbf{Stable within-dataset identifiers}: Records in each source carry a \localentity{} identifier that is consistent for all records belonging to the same \globalentity{} \textit{within that dataset}.
	      Between datasets, the identifier schema differs ($DF1_*$ vs $DF2_*$), but the within-dataset stability assumption scopes us away from clustering/graph-based methods that would otherwise be necessary to discover local groups.
	\item \textbf{Attribute types}: Matching signals derive from textual, numerical, and datetime columns only—scoping out LLM-based semantic matching and behavioural/network-based methods.
	\item \textbf{Cost minimisation}: The objective is to minimise total resolution error (or equivalently, maximise $F_1$) under computational budget constraints imposed by blocking.
\end{enumerate}

These constraints narrow the methodological space substantially: the \fellegisunter{} probabilistic model with expectation-maximisation (\emalgo{}) estimation, implemented via \splink{}, emerges as the natural (and arguably \textbf{only}) principled choice for matching.

\subsection{Contributions}
\label{sec:contributions}

\begin{enumerate}
	\item \textbf{Formalised Local-to-Global Framework} (\Cref{sec:framework}): A clear separation between local probabilistic matching (\fellegisunter{} via \splink{}) and global aggregation, with explicit assumptions and failure modes.
	\item \textbf{Implementation Patterns for Unsupervised \splink{}} (\Cref{sec:implementation}): Blocking strategy design with recall bounds, custom comparison levels (exact, \jw{}, percentage-difference, temporal), \emalgo{} training schedules, and evaluation without labels.
	\item \textbf{Global Aggregation Taxonomy} (\Cref{sec:agg_functions}): Mathematical unification of naive sum, expected-score normalisation, noisy-OR, max-over-partitions, and vector-space aggregation.
	\item \textbf{Experimental Validation} (\Cref{sec:results}): Systematic results on synthetic data and public datasets (MusicBrainz).
	\item \textbf{Cardinality Extremities \& Layer-2 Mitigation} (\Cref{sec:edge_cases}): Diagnosis of independence-assumption breakdown at low cardinality; proposal for alternative approaches.
\end{enumerate}

\subsection{Report Structure}
\label{sec:structure}

This report alternates focus between \textit{rationale} (why a design decision was made) and \textit{practical implementation} (how it was realised in code).
\begin{itemize}
	\item \textbf{Section 2}: Theoretical framework—\fellegisunter{} foundations, \splink{} implementation, global aggregation framework, hierarchical probabilistic linkage.
	\item \textbf{Section 3}: System implementation—data generation, \splink{} configuration, blocking strategy.
	\item \textbf{Section 4}: Experimental methodology—datasets, metrics, evaluation protocol.
	\item \textbf{Section 5}: Experimental results—artificial baseline, public dataset validation, behavioural aggregation.
	\item \textbf{Section 6}: Edge cases—cardinality failure modes.
	\item \textbf{Section 7}: Discussion—confidence scores necessity, practical takeaways.
	\item \textbf{Section 8}: Future work and conclusion.
	\item \textbf{Annexes}: ER methodology summary, mathematical derivations
\end{itemize}
